\begin{lemma}[Right presentation invariance]
If $\mathsf{PowerPresentationEq}(Y,Y')$, then for every
$X\in\mathsf{PowerQ}(Q)$,
\[
\mathsf{PowerRel}(r)(X,Y)\quad\Longleftrightarrow\quad
\mathsf{PowerRel}(r)(X,Y').
\]
\end{lemma}
\begin{proof}
Fix $r$ and, for a presentation $Y$, let $P(Y)$ be the assertion that for
every $Y'$ and every $X$, structural equivalence
$\mathsf{PowerPresentationEq}(Y,Y')$ implies
\[
\mathsf{PowerRel}(r)(X,Y)\quad\Longleftrightarrow\quad
\mathsf{PowerRel}(r)(X,Y').
\]
We prove $P(Y)$ by well-founded induction on the child-to-parent relation
\noderef{PowerChild}, whose well-foundedness is
\noderef{PowerChildWellFounded}.  At every use of the induction hypothesis
below, the first right presentation is a child of the current right
presentation: if $Y=\mathsf{node}(\kappa,h,g)$, then
\[
\mathsf{PowerChild}(g(j),Y)
\]
is witnessed, according to the definition in \noderef{PowerChild}, by
the pair $(j,\mathsf{rfl})$.  The four reductions used in the cases are the
corresponding branches of the well-founded recursion defining
\noderef{PowerRel}.

First suppose that $Y=\mathsf{atom}(q)$ and
$Y'=\mathsf{atom}(q')$.  The atom--atom clause of
\noderef{PowerPresentationEq} gives $q=q'$.  If
$X=\mathsf{atom}(p)$, both sides reduce by
\noderef{PowerRelAtomAtom} to $r(p,q)$ after substituting this equality.
If $X=\mathsf{node}(\iota,h,f)$, both sides reduce by
\noderef{PowerRelNodeAtom} to the universal comparison of every $f(i)$
with the same atom, again because $q=q'$.  Thus the equivalence holds for
both possible left tags.

The two mixed right-tag cases cannot occur: if one of $Y,Y'$ is an atom
and the other is a node, the corresponding branch of
\noderef{PowerPresentationEq} is false.  Hence the assumed structural
equivalence rules out both cases.

It remains to consider
\[
Y=\mathsf{node}(\kappa,h,g),\qquad
Y'=\mathsf{node}(\lambda,h',g').
\]
Write the two conjuncts of the assumed equivalence as
\[
\begin{split}
 (E_\to)&:\ \forall j\in\kappa\;\exists j'\in\lambda,
   \mathsf{PowerPresentationEq}(g(j),g'(j')),\\
 (E_\leftarrow)&:\ \forall j'\in\lambda\;\exists j\in\kappa,
   \mathsf{PowerPresentationEq}(g(j),g'(j')).
\end{split}
\]
These are exactly the forward and backward matching clauses in
\noderef{PowerPresentationEq}.  Equivalently, by the symmetry part of
\noderef{PowerPresentationEqEquivalence}, the second clause is the forward
matching clause for the symmetric equivalence
$\mathsf{PowerPresentationEq}(Y',Y)$; this permits the same argument to be
read in either direction.

Suppose first that $X=\mathsf{atom}(p)$.  By
\noderef{PowerRelAtomNode}, the left side is equivalent to
\[
\exists j\in\kappa\quad
\mathsf{PowerRel}(r)(\mathsf{atom}(p),g(j)).
\]
Given such a $j$, choose $j'$ from $(E_\to)$.  The induction hypothesis
applies to the matched pair $g(j),g'(j')$, because the explicit witness
$(j,\mathsf{rfl})$ proves $\mathsf{PowerChild}(g(j),Y)$.  It changes the
displayed comparison into
\[
\mathsf{PowerRel}(r)(\mathsf{atom}(p),g'(j')).
\]
The existential clause of \noderef{PowerRelAtomNode} then proves the right
side.  Conversely, start with a witness $j'$ for the right side, choose
$j$ from $(E_\leftarrow)$, and apply the same induction hypothesis, now
using the child witness $(j,\mathsf{rfl})$ for $Y$.  This time use the
reverse implication supplied by the induction equivalence, and then use the
atom--node reduction to return to the left side.

Now suppose that $X=\mathsf{node}(\iota,k,f)$.  By
\noderef{PowerRelNodeNode}, the left side is equivalent to
\[
\forall i\in\iota\;\exists j\in\kappa\quad
\mathsf{PowerRel}(r)(f(i),g(j)).
\]
Fix an arbitrary $i$ and choose a corresponding $j$.  Use $(E_\to)$ to
choose $j'$ with
$\mathsf{PowerPresentationEq}(g(j),g'(j'))$.  The induction hypothesis,
justified by $\mathsf{PowerChild}(g(j),Y)$ via $(j,\mathsf{rfl})$, gives
\[
\mathsf{PowerRel}(r)(f(i),g'(j')).
\]
Thus the
same arbitrary left-child index $i$ has a right-child witness $j'$ in
$Y'$.  Applying \noderef{PowerRelNodeNode} proves the right side.  For the
converse, begin with an arbitrary $i$ and a witness $j'$ supplied by the
right-hand node--node reduction, use $(E_\leftarrow)$ to obtain $j$, and
apply the reverse implication of the induction hypothesis along the child
witness $(j,\mathsf{rfl})$.  This supplies the required $j$ for the
left-hand node--node reduction.  In particular, the universal
quantification over the left children is preserved; only its existential
right-child witness is transported.

The two implications have therefore been proved for both possible left
tags, while the atom and mixed right-tag cases were already settled.  This
establishes $P(Y)$ for every $Y$, and specializing it to the given
$X,Y,Y'$ and $hyy'$ yields the claimed right presentation invariance.
\end{proof}
